\hypertarget{InitialConditionsFileFormat}{}
\subsection{initial\_conditions\_map.txt\label{sec_initialConditionsMap}}
\index{Initial conditions}
%\hyperlink{InitialConditionsFileFormat}{Goto InitialConditionsFileFormat}

%
\subsubsection{Purpose}
The well-posedness of the shallow water problem requires initial conditions and boundary conditions be provided. The present section deals with the specification of initial conditions.

The initial conditions may be either uniform or spatially distributed. They are provided on a cell-by-cell basis. Two types of variable must be provided: the water depth or free surface elevation, and the flow velocity or unit discharge components.

%
\subsubsection{Format description}
%
File format for a uniform initial condition.  \texttt{\textbf{Bold}} : keyword, must appear as such in the file. \texttt{<Tab>}: tabulation character. \\
    \begin{tabular}{p{6cm} | p{9cm}}
    \hline
    Field(s) & Comment(s) \\
    \hline  \hline
    \texttt{\textbf{Unif} <Tab> \textbf{1}} & Flag is 1 because the initial conditions are uniform \\ \hline
    \texttt{\textbf{==== Default Param}} & Section separator line, leave unchanged \\  \hline
    \texttt{\textbf{z}} or \texttt{\textbf{h}} \texttt{<Tab>} \texttt{\textbf{u}} or \texttt{\textbf{q}} \texttt{<Tab>} \texttt{\textbf{v}} or \texttt{\textbf{r}} & \texttt{\textbf{h}}: the first numerical value in the record is the water depth \\ 
     & \texttt{\textbf{z}}: the first numerical value in the record is the surface elevation \\
     & \texttt{\textbf{u}}: the second numerical value in the record is the $x-$ flow velocity\\ 
     & \texttt{\textbf{q}}: the second numerical value in the record is the $x-$ unit discharge\\ 
     & \texttt{\textbf{v}}: the third numerical value in the record is the $y-$ flow velocity\\ 
     & \texttt{\textbf{r}}: the third numerical value in the record is the $y-$ unit discharge\\ 
    \hline
     \texttt{z} or \texttt{h} \texttt{<Tab>} \texttt{u} or \texttt{q} \texttt{<Tab>} \texttt{v} or \texttt{r} & The three numerical values for the uniform initial condition \\      
     \texttt{\textbf{==== Distrib}} & Section separator line, leave unchanged \\  
     \hline
    \end{tabular}
%
\vspace{10pt}

File format for a spatially variable initial condition.  \texttt{\textbf{Bold}} : keyword, must appear as such in the file. \texttt{<Tab>}: tabulation character. \\
    \begin{tabular}{p{6cm} | p{9cm}}
    \hline
    Field(s) & Comment(s) \\
    \hline  \hline
    \texttt{\textbf{Unif} <Tab> \textbf{1}} & Flag is 0 because the initial conditions are spatially variable \\ \hline
    \texttt{\textbf{==== Default Param}} & Section separator line, leave unchanged \\  \hline
    \texttt{Val\_1 <Tab> Val\_2 <Tab> Val\_3}  & These 3 numerical values will be ignored because \texttt{Flag} = 0 \\ \hline
    \texttt{\textbf{==== Distrib}} & Section separator line, leave unchanged \\  \hline
    \texttt{\textbf{z}[i]} or \texttt{\textbf{h}[i]} \texttt{<Tab>} \texttt{\textbf{u}[i]} or \texttt{\textbf{q}[i]} \texttt{<Tab>} \texttt{\textbf{v}[i]} or \texttt{\textbf{r}[i]} 
    & \texttt{\textbf{h}}: the first numerical value in the record is the water depth \\ 
     & \texttt{\textbf{z}}: the first numerical value in the record is the surface elevation \\
     & \texttt{\textbf{u}}: the second numerical value in the record is the $x-$ flow velocity\\ 
     & \texttt{\textbf{q}}: the second numerical value in the record is the $x-$ unit discharge\\ 
     & \texttt{\textbf{v}}: the third numerical value in the record is the $y-$ flow velocity\\ 
     & \texttt{\textbf{r}}: the third numerical value in the record is the $y-$ unit discharge\\ 
     \hline
     \texttt{z[i]} or \texttt{h[i]} \texttt{<Tab>} \texttt{u[i]} or \texttt{q[i]} \texttt{<Tab>} \texttt{v[i]} or \texttt{r[i]} 
    &Initial conditions for cell $i$, index $i$ is the index of the cell in the mesh \\
    \hline
     \end{tabular}
%
\subsubsection{Sample files}
%
\paragraph*{Example \#1: uniform initial conditions (elevation-velocity)}
The initial condition is uniform. The free surface elevation (1 m) and the flow velocity (0 m/s in both directions of space) are the same in all cells.
\begin{verbatim}
Unif	1
==== Default Param
z	u	v
1	0	0
==== Distrib
\end{verbatim}
%
\paragraph*{Example \#2: uniform initial conditions (depth-unit discharge)}
The initial condition is uniform. The water depth (2 m) and the unit discharge (1 m$^2$/s in the $x-$direction, 0 m$^2$/s in the $y-$direction) are the same in all cells.
\begin{verbatim}
Unif	1
==== Default Param
h	q	r
2	1	0
==== Distrib
\end{verbatim}
%
\paragraph*{Example \#3: distributed initial conditions}
There are 3 cells in the model. In Cell\#1, the initial free surface elevation is 10m, the flow velocity is zero. In Cell\#2, the free surface elevation is 10m, the $x-$velocity is 1 m/s and the $y-$velocity is zero. In the third cell, the free surface elevation is 5 m, and the initial flow velocity is zero. 
\begin{verbatim}
Unif	0
==== Default Param
z	u	v
1	0	0
==== Distrib
10. 0.  0.
10. 1.  0.
5.  0.  0.
\end{verbatim}